\documentclass[tikz,border=2bp]{standalone}
\usepackage{ctex}
\usepackage[europeanresistors, americanvoltages, americancurrents, siunitx]{circuitikz}
\standaloneenv{circuitikz}

\begin{document}

  % ==================== (a) 原电路 ====================
  \begin{circuitikz}[semithick, scale=1.0, transform shape]
    % 5V电源
    \draw (0,0) to[vsource, l={\SI{5}{\volt}}] (0,3)
          to[R, l={$25\text{ k}\Omega$}] (2.5,3);
          
    % 双掷开关触点与闸刀
    \node[circle, fill, inner sep=1.2pt] (P1) at (2.5, 3) {};
    \node[above] at (2.5,3) {1};
    \node[circle, fill, inner sep=1.2pt] (P2) at (2.5, 1.2) {};
    \node[left] at (2.3,1.2) {2};
    \node[circle, fill, inner sep=1.2pt] (P0) at (4.0, 3) {};
    
    % 开关闸刀 (指向1，带转向2的箭头)
    \draw[very thick] (4.0,3) -- (2.6,3.1);
    \draw[->, >=stealth, thin] (2.8, 2.9) arc (160:200:1.8);
    \node[left] at (3.4, 2.3) {$t=0$};
    
    % 连接位置2的电阻
    \draw (2.5,1.2) to[R, l={$100\text{ k}\Omega$}] (2.5,0);
    
    % 动端连接右侧 C 和 R1 的并联
    \draw (4.0,3) -- (5.0,3)
          to[C, l_={$10\text{ }\mu\text{F}$}] (5.0,0);
    \draw (5.0,3) -- (6.5,3)
          to[R, l={$100\text{ k}\Omega$}] (6.5,0);
          
    % 底部地线
    \draw (6.5,0) -- (0,0);
    \draw (3.25,0) node[ground] {};
    
    \node[below] at (3.25,-0.5) {(a) 原电路};
  \end{circuitikz}

  % ==================== (b) t=0- 稳态等效电路 ====================
  \begin{circuitikz}[semithick, scale=1.0, transform shape]
    % 5V电源与25k电阻
    \draw (0,0) to[vsource, l={\SI{5}{\volt}}] (0,3)
          to[R, l={$25\text{ k}\Omega$}] (2.5,3)
          -- (4.0,3);
          
    % 电容在稳态下开路 (极性及名称标注在右侧)
    \draw (4.0,3) -- (5.0,3)
          to[open] (5.0,0);
    \node[right] at (5.2,2.2) {$+$};
    \node[right] at (5.2,0.8) {$-$};
    \node[right] at (5.2,1.5) {$u_C(0^-)$};
    
    % 100k电阻R1
    \draw (5.0,3) -- (6.5,3)
          to[R, l={$100\text{ k}\Omega$}] (6.5,0);
          
    % 底部地线
    \draw (6.5,0) -- (0,0);
    \draw (3.25,0) node[ground] {};
    
    % 结果文本框 (移动至 x=2.2 避开左侧电源与右侧极性)
    \node[draw, align=left, font=\small] at (2.2,1.1) {
      $u_C(0^-) = \SI{4}{\volt}$
    };
    
    \node[below] at (3.25,-0.5) {(b) $t=0^-$ 稳态等效电路};
  \end{circuitikz}

  % ==================== (c) t>=0 放电电路 ====================
  \begin{circuitikz}[semithick, scale=1.0, transform shape]
    % 电容在最左侧
    \draw (0,3) to[C, l_={$10\text{ }\mu\text{F}$}] (0,0);
    \node[left] at (-0.2,1.5) {$u_C(t)$};
    \node[left] at (-0.2,2.2) {$+$};
    \node[left] at (-0.2,0.8) {$-$};
    % 电容电流标注 (放电电流向下)
    \node[right] at (0.2,2.2) {$i_C(t)$};
    \draw[->, >=stealth, thin] (0.2,2.5) -- (0.2,1.9);
    
    % 并联的电阻R2 (左侧，即开关拨到的那个电阻，流过电流i)
    \draw (0,3) -- (2.0,3)
          to[R, l_={$100\text{ k}\Omega$}, i_>={$i(t)$}] (2.0,0);
          
    % 并联的电阻R1 (右侧)
    \draw (2.0,3) -- (4.0,3)
          to[R, l={$100\text{ k}\Omega$}] (4.0,0);
          
    % 底部地线
    \draw (4.0,0) -- (0,0);
    \draw (2.0,0) node[ground] {};
    
    \node[below] at (2.0,-0.5) {(c) $t\ge0$ 放电电路};
  \end{circuitikz}

\end{document}
